\documentclass[ProjectKWK]{subfiles}
% WARNING: AuCTeX local variables only get reset when file is loaded
% and differ between this file and ProjectKWK.tex
% so must re-load whichever file you want to compile with C-x C-v

% WARNING: Different AucTeX execution depending on whether
% 0. Being compiled as standalone document
%    * Compile main once
%    * Then compile this one
%    * Keep compiling until nothing changes
% 0. Being compiled as subfile of main document
%    * Just compile main document repeatedly

\providecommand{\econtexRoot}{}
\renewcommand{\econtexRoot}{..}
\providecommand{\econtexPaths}{}\renewcommand{\econtexPaths}{\econtexRoot/Resources/econtexPaths}
% The \commands below are required to allow sharing of the same base code via Github between TeXLive on a local machine and Overleaf (which is a proxy for "a standard distribution of LaTeX").  This is an ugly solution to the requirement that custom LaTeX packages be accessible, and that Overleaf prohibits symbolic links
\providecommand{\econtex}{\econtexRoot/Resources/texmf-local/tex/latex/econtex}
\providecommand{\econtexSetup}{\econtexRoot/Resources/texmf-local/tex/latex/econtexSetup}
\providecommand{\econtexShortcuts}{\econtexRoot/Resources/texmf-local/tex/latex/econtexShortcuts}
\providecommand{\econtexBibMake}{\econtexRoot/Resources/texmf-local/tex/latex/econtexBibMake}
\providecommand{\econtexBibStyle}{\econtexRoot/Resources/texmf-local/bibtex/bst/econtex}
\providecommand{\econtexBib}{\econtexRoot/Resources/texmf-local/bibtex/bib/economics}
\providecommand{\notes}{\econtexRoot/Resources/texmf-local/tex/latex/handout}
\providecommand{\handoutSetup}{\econtexRoot/Resources/texmf-local/tex/latex/handoutSetup}
\providecommand{\handoutShortcuts}{\econtexRoot/Resources/texmf-local/tex/latex/handoutShortcuts}
\providecommand{\handoutBibMake}{\econtexRoot/Resources/texmf-local/tex/latex/handoutBibMake}
\providecommand{\handoutBibStyle}{\econtexRoot/Resources/texmf-local/bibtex/bst/handout}

\providecommand{\FigDir}{\econtexRoot/Figures}
\providecommand{\CodeDir}{\econtexRoot/Code}
\providecommand{\DataDir}{\econtexRoot/Data}
\providecommand{\SlideDir}{\econtexRoot/Slides}
\providecommand{\TableDir}{\econtexRoot/Tables}
\providecommand{\ApndxDir}{\econtexRoot/Appendices}

\providecommand{\ResourcesDir}{\econtexRoot/Resources}
\providecommand{\rootFromOut}{..} % Path back to root directory from output-directory
\providecommand{\LaTeXGenerated}{\econtexRoot/LaTeX} % Put generated files in subdirectory
\providecommand{\econtexPaths}{\econtexRoot/Resources/econtexPaths}
\providecommand{\LaTeXInputs}{\econtexRoot/Resources/LaTeXInputs}
\providecommand{\LtxDir}{LaTeX/}
\providecommand{\EqDir}{Equations} % Put generated files in subdirectory

\onlyinsubfile{% https://tex.stackexchange.com/questions/463699/proper-reference-numbers-with-subfiles
    \csname @ifpackageloaded\endcsname{xr-hyper}{%
      \externaldocument{\econtexRoot/BufferStockTheory}% xr-hyper in use; optional argument for url of main.pdf for hyperlinks
    }{%
      \externaldocument{\econtexRoot/BufferStockTheory}% xr in use
    }%
    \renewcommand\labelprefix{}%
    % Initialize the counters via the labels belonging to the main document:
    \setcounter{equation}{\numexpr\getrefnumber{\labelprefix eq:Dummy}\relax}% eq:Dummy is the last number used for an equation in the main text; start counting up from there
}


\onlyinsubfile{\externaldocument{ProjectKWK}} % Get xrefs -- esp to appendix -- from main file; only works properly if main file has already been compiled;

\begin{document}

% Attempted to make all lines used for Web version contain {Web} (or version with only single curly brace at end) so can be removed with sed
\providecommand{\versn}{pdf} % Version; like, web or pdf or journal submission
\ifthenelse{\boolean{Web}}{    % {Web}
  \renewcommand{\versn}{Web}     % Too hard to figure out passing -output-directory through make4ht through htlatex, so web version is compiled with junk files in main directory
  \renewcommand{\rootFromOut}{.} % {Web}
}{}  % {Web}

% Tiny info header at top to track git commit
%\hfill{\tiny \jobname~\versn~\today~{at} \DTMcurrenttime, \input{\ResourcesDir/.git-source-commit}~~\input{\ResourcesDir/.git-public-commit}}

\title{Close Elections, Campaign Contributions,\\ and Financial Deregulation}

\author{Kyung Woong Koh\authNum}

\keywords{Financial Deregulation, Close Elections}

\jelclass{D81, D91, E21 \par
  \href{https://econ-ark.org}{\includegraphics{\ResourcesDir/PoweredByEconARK}}
}

\renewcommand{\forcedate}{November 9, 2021}\date{\forcedate}
\maketitle
\hypertarget{abstract}{}
\begin{abstract}
This paper builds upon Igan and Mishra (2014) \cite{IM14} on vote switches towards financial deregulation by US legislators. I measure the effect of close elections on US legislators on switching their votes towards financial deregulation in Congress bills. I aim to distinguish between vote switches towards financial deregulation because of voters' general interests (especially after the Global Financial Crisis of 2007) versus the financial industry's special interests and the industry's campaign contributions and lobbying expenditures towards legislators in close elections.
\end{abstract}


\begin{authorsinfo}
  \name{Contact: \href{mailto:kkoh8@jhu.edu}{\texttt{kkoh8@jhu.edu}}, Department of Economics, Johns Hopkins University, Baltimore, MD 21218.}
\end{authorsinfo}

\newcommand{\thankstext}{
  The paper is the author's 2nd year paper for fulfillment to the PhD in Economics program at Johns Hopkins University, and also for Professor Christpher Carroll's Computational Methods course. I thank Professors Laurence Ball and Filipe Campante for comments and advice, and Deniz Igan and Prachi Mishra for sharing their data.}

\ifthenelse{\boolean{Web}}{}{
  \begin{minipage}{0.9\textwidth}
    \tiny \thankstext
\end{minipage}
} % {Web}
{\titlepagefinish}

\hypertarget{Introduction}{}
\section{Introduction}\label{sec:intro}


\begin{figure}[ht]
  \centerline{
    \includegraphics[width=0.8\textwidth]{\FigDir/KohFigure1_tikzmake}
  }
    \caption{Two Channels in which Legislators in Close Elections Switch their Votes Toward Financial Deregulation}\label{KohFigure1}
\end{figure}


How do legislators in close elections vote differently on financial deregulation compared to legislators in sure elections? I use data from \cite{IM14} to answer this question.


\hypertarget{The-Problem}{}

\section{The Problem}

Legislators, to get reelected, have to respond to voters and special interests (in this paper, the financial industry)

Legislators in sure elections are less incentivized to switch their vote on bills on financial deregulation; they will still win their sure elections whether they vote for (if the legislator's constituent voters prefer deregulation) or against (if the legislator's constituent voters do not prefer deregulation). Legislators in sure elections are also less incentivized to switch their votes towards financial deregulation to satisfy the financial industry, as they will still have enough campaign finances to win their sure elections regardless of the additional campaign contributions from the financial industry. In short, legislators in sure elections would not be strongly incentivized to switch their votes on financial deregulation to get more votes or acquire more campaign contributions. Legislators in sure elections will be less responsive to additional lobbying expenditures or campaign contributions by the financial industry as well.

Legislators in close elections may be more incentivized to switch their vote towards financial deregulation via the same two stakeholders. To attract more votes, legislators in close elections may be more responsive to voters' needs for financail deregulation. To gather more campaign contributions, those legislators may also be more responsive to lobbying expenditures and campaign contributions from the financial industry. Legislators in close elections, through either mechanism, have a greater incentive to switch their votes towards financial deregulation in order to win reelection.

\cite{IM14} argue that their results, that lobbying expenditures make legislators more likely to switch their votes towards financial deregulation, is a correlational result and therefore cannot be taken as causation from either of these reasons. I attempt to provide at least a partial answer by introducing new variables, including one for whether a legislator has been in or will be facing a close election. I also include other variables to support the regression strategies that can differentiate between voters' public scrutiny and the financial industry's special interests.

There are two broad strategies that can be used; the first is to use a measure of ``media congruence'', by Snyder and Stromberg (2010).

The second is to use the rapid change in public sentiment against financial deregulation shortly after the Global Financial Crisis and the resulting Great Recession.





\subsection{Setup}\label{subsec:Setup}





\hypertarget{Estimation Strategy}{}

\section{Variables}

\subsection{Dependent Variable}\label{subsec:DependentVariable}

\newsavebox{\DepVar}
\begin{table}
\centering
\caption{Definition of the Main Dependent Variable, Vote Switch towards Deregulation}\label{table:DepVar}
\sbox{\DepVar}{
\begin{tabular}{p{0.25\linewidth} | p{0.25\linewidth} | p{0.25\linewidth}} \hline
\textbf{Value of $S_iBR$}  & Voted for deregulation in Bill $B,R$ & Voted against deregulation in Bill $B,R$ \\ \hline
Voted for deregulation in Bill $B,R-1$ & 0 & 0 \\ \hline
Voted for deregulation in Bill $B,R-1$ & 1 & 0 \\ \hline
\end{tabular}
}
\usebox{\DepVar}
\end{table}

\subsection{Election Closeness}\label{subsec:ElectionCloseness}

Compared to the original \cite{IM14} paper, I add the new variable of ``Election Closeness'' as the main focus of this paper. I define ``election closeness'' for each legislator $i$ and bill $BR$, denoted as $X_{iBR}$, as the degree to which the legislator has faced (past) or will face (future) a close election. Note the two possible ways of looking at election closeness: through the past election(s) of the legislator, or to the (immediately next) future election of the legislator. Meanwhile, the measure of future election closeness
\\

Past election closeness is defined simply, in this paper, as the percentage margin of victory of the legislator in the last election before a vote on a bill $BR$. I denote this variable as $X^{P}_{iBR}$. As a concrete example, assume that legislator $A$ won his/her last election with 49,000 votes against the runner-up, who got 47,000 votes in a congressional election of a total of 100,000 votes, and that the remaining 4,000 votes all went to third-party, independent, and write-in votes. In this case, legislator $A$'s margin of victory is $(49,000-47,000)/100,000 = 2\%$. There are two important characteristics of this variable: first, $X^{P}_{iBR}$ must necessarily be greater than zero, as the legislator must have won at least one more vote than the runner-up. Second, $X^{P}_{iBR}$ is the same across all bills $BR$ during the same congress $C$, as legislator $i$ has only one value of past election closeness at any given congress $C$ and all the bills therein. \\

Ideally, future election closeness may be defined as the expected margin of victory in the next (future) election of the legislator. I denote this variable as $X^{F}_{iBR}$. The future expected margin of victory of an as-of-yet undecided election can be proxied by results of election polls. Since most future congressional elections at any given electoral cycle have at least one polling result, and in most cases more than one, $X^{F}_{iBR}$ can differ across bills $BR$ even in the same congress $C$.


\subsection{Original Variables}\label{subsec:OriginalVariables}



\hypertarget{Regressions}{}

\section{Regressions}

\subsection{Regression A}\label{subsec:Regression A}


Concretely, I write Regression A1 as:

\begin{align}
    S_{iBR} &= \beta_{1} L_{BR} + \beta_{2} X_{iBR}^{P} + \beta_{3} (L_{BR} \times X_{iBR}^{P} ) \nonumber \\
    &+ \alpha F_{BR} + \gamma T_{BR}+ s_{i} \times t_{c}+ v_{B} \times t_{c}+\mu_{R} \times t_{c}+\varepsilon_{iBR}
\end{align}



\hypertarget{Conclusions}{}
\section{Conclusions}


% %The paper's results are all easily reproducible \href{https://econ-ark.org/_materials/BufferStockTheory?launch}{interactively on the web} or \href{https://github.com/econ-ark/BufferStockTheory}{on any standard computer system}.  Such reproducibility reflects the paper's use of the open-source \href{https://econ-ark.org}{Econ-ARK} toolkit, which is used to generate all of the quantitative results of the paper, and which integrally incorporates all of the analytical insights of the paper.

% % The Dummy equation below sems to be needed to get the equation numbering in the appendix
% % reliably to start at the next number after the last actual equation number used in the paper

% \clearpage\vfill\eject
% \begin{equation*}
%   \label{eq:Dummy}
% \end{equation*}

% \onlyinsubfile{\bibliography{
%     \texname, % subfile inherits texname from preamble of parent
%     \econtexBib % Default bib database is in Resources/LaTeXInputs
%   }}

\onlyinsubfile{% Allows two (optional) supplements to hard-wired \texname.bib bibfile:
% economics.bib is a default bibfile that supplies anything missing elsewhere
% Add-Refs.bib is an override bibfile that supplants anything in \texfile.bib or economics.bib
\IfFileExists{\econtexRoot/Add-Refs.bib}{
  % then 
  \typeout{References in Add-Refs.bib will take precedence over those elsewhere}
  \setboolean{AddRefsExists}{true}
  \setboolean{NeitherExists}{false} % Default is true
}{
  % else
  \setboolean{AddRefsExists}{false} % No added refs exist so defaults will be used
  \setboolean{BothExist}{false}     % Default is that Add-Refs and economics.bib both exist
}

% Deal with case where economics.bib is found by kpsewhich 
\IfFileExists{/usr/local/texlive/texmf-local/bibtex/bib/economics.bib}{
  % then
  \typeout{References in default global economics.bib will be used for items not found elsewhere}
  \setboolean{economicsExists}{true}
  \setboolean{NeitherExists}{false}
}{
  % else 
  \typeout{Found no global database file}
  \setboolean{economicsExists}{false}
  \setboolean{BothExist}{false}
}

\IfFileExists{economics.bib}{
  % then
  \typeout{References in economics.bib will be used for items not found elsewhere}
  \setboolean{economicsExists}{true}
  \setboolean{NeitherExists}{false}
}{
  % else 
  \typeout{Found no global database file}
  \setboolean{economicsExists}{false}
  \setboolean{BothExist}{false}
}

\ifthenelse{\boolean{BothExist}}{
  % then use both
  \typeout{bibliography{\econtexRoot/Add-Refs,\econtexRoot/\texname,economics}}
  \bibliography{\econtexRoot/Add-Refs,\econtexRoot/\texname,economics}
  % else both do not exist
}{ % maybe neither does?
  \ifthenelse{\boolean{NeitherExists}}{
    \typeout{bibliography{\texname}}
    \bibliography{\texname}}{
    % no -- at least one exists
    \ifthenelse{\boolean{AddRefsExists}}{
      \typeout{bibliography{\econtexRoot/Add-Refs,\econtexRoot/\texname}}
      \bibliography{\econtexRoot/Add-Refs,\econtexRoot/\texname}}{
      \typeout{bibliography{\econtexRoot/\texname,economics}}
      \bibliography{         \econtexRoot/\texname,economics}}
  } % end of picking the one that exists
} % end of testing whether neither exists
}
%\bibliography{economics}
\end{document}
\endinput

% If you are editing in Emacs-AucTeX, modify the lines below for your system (otherwise ignore)
% Local Variables:
% eval: (setq TeX-command-list  (assq-delete-all (car (assoc "BibTeX" TeX-command-list)) TeX-command-list))
% eval: (setq TeX-command-list  (assq-delete-all (car (assoc "BibTeX" TeX-command-list)) TeX-command-list))
% eval: (setq TeX-command-list  (assq-delete-all (car (assoc "BibTeX" TeX-command-list)) TeX-command-list))
% eval: (setq TeX-command-list  (assq-delete-all (car (assoc "Biber"  TeX-command-list)) TeX-command-list))
% eval: (add-to-list 'TeX-command-list '("BibTeX" "bibtex LaTeX/%s" TeX-run-BibTeX nil t                                                                              :help "Run BibTeX") t)
% eval: (add-to-list 'TeX-command-list '("BibTeX" "bibtex LaTeX/%s" TeX-run-BibTeX nil (plain-tex-mode latex-mode doctex-mode ams-tex-mode texinfo-mode context-mode) :help "Run BibTeX") t)
% TeX-PDF-mode: t
% TeX-file-line-error: t
% TeX-debug-warnings: t
% LaTeX-command-style: (("" "%(PDF)%(latex) %(file-line-error) %(extraopts) -output-directory=LaTeX %S%(PDFout)"))
% TeX-source-correlate-mode: t
% TeX-parse-self: t
% eval: (cond ((string-equal system-type "darwin")    (progn (setq TeX-view-program-list '(("Skim" "/Applications/Skim.app/Contents/SharedSupport/displayline -b %n LaTeX/%o %b"))))))
% eval: (cond ((string-equal system-type "gnu/linux") (progn (setq TeX-view-program-list '(("Evince" "evince --page-index=%(outpage) LaTeX/%o"))))))
% eval: (cond ((string-equal system-type "gnu/linux") (progn (setq TeX-view-program-selection '((output-pdf "Evince"))))))
% TeX-parse-all-errors: t
% End:
